% -*- coding: utf-8; time-stamp-format: "%02d-%02m-%:y at %02H:%02M:%02S %Z" -*-
%<*none>
\def\dtxtimestamp {Time-stamp: <03-10-2016 at 23:00:16 CEST>}
%</none>
%<*!readme>
%%
%% Package: footnotehyper
%% Version: 0.9f (2016/10/03)
%% License: LPPL 1.3c
%% Copyright (C) 2016 Jean-Francois B. <jfbu on github>.
%%
%</!readme>
%<*insfile|tex>
\def\pkgname        {footnotehyper}
\def\pkgdate        {2016/10/03}
\def\docdate        {2016/10/03}
\def\pkgversion     {v0.9f}
\def\pkgdescription {hyperref aware footnote.sty (JFB)}
%</insfile|tex>
%<*none>
% Definition of \pkgLicense
\begingroup% cette méthode ne marcherait pas avec caractères en dehors de 32-127
    \long\def\firstofone #1{#1}\catcode1=14\catcode2=0
    \catcode`\%=12\catcode`\_=12\endlinechar13\catcode13=13 ^^A
    \catcode32=13\catcode`\\=12^^Brelax^^A
^^Bfirstofone{^^Bendgroup^^Bdef^^BpkgLicense^^A
{% Package: footnotehyper
% Version: 0.9f (2016/10/03)
% License: LPPL 1.3c
% Copyright (C) 2016 Jean-Francois B. <jfbu on github>.
%
% This Work may be distributed and/or modified under the conditions
% of the LaTeX Project Public License, version 1.3c. This version of
% this license is in:
%
% > <http://www.latex-project.org/lppl/lppl-1-3c.txt>
%
% and the latest version of this license is in:
%
% > <http://www.latex-project.org/lppl.txt>
%
% Version 1.3 or later is part of all distributions of
% LaTeX version 2005/12/01 or later.
%
% The Author of this Work is: Jean-Francois B. `<jfbu on github>`
%
% This Work consists of the main source file footnotehyper.dtx and the
% derived files footnotehyper.sty, footnotehyper.ins, footnotehyper.tex,
% footnotehyper.pdf, footnotehyper.dvi.
}}%
\begingroup\catcode1 0 \catcode`\\ 12 
^^Aiffalse
%</none>
%<*readme>
<!-- -->

    Source:  footnotehyper.dtx (v0.9f 2016/10/03)
    Author:  Jean-Francois B.
    Info:    hyperref aware footnote.sty
    License: LPPL 1.3c
    Copyright (C) 2016 Jean-Francois B. <jfbu on github>.


ABSTRACT
========

The `footnote` package by Mark Wooding (which dates back to 1997)
allows to gather footnotes (`\savenotes`) and later insert them
(`\spewnotes`) at the bottom of the page, even if the intervening
material consists of tabulars, minipages or framed contents for
example. There is a command to patch any given environment to do
the `\savenotes/\spewnotes` commands automatically. Also there is
a `footnote` environment which allows to insert verbatim material.

The `footnotehyper` package loads `footnote.sty` and rewrites a
number of its commands to add `hyperref` awareness and to fix some
issues, particularly the one of incompatibility with
`color/xcolor`. The incompatibility with `babel-frenchb` is
addressed, too. In case of absence of `hyperref`, `footnotehyper`
limits itself to partially fix some of these issues.

This package provides no facility dedicated to footnotes from
floating environments.


INSTALLATION
============

To extract the package (.sty) and driver (.tex) files from
footnotehyper.dtx:

- if footnotehyper.ins is present:   etex footnotehyper.ins
- without footnotehyper.ins:         etex footnotehyper.dtx
- or run latex or pdflatex directly on footnotehyper.dtx

At least three ways to produce footnotehyper.pdf (method (1) is
preferred):

1. latex footnotehyper.tex (twice), then dvipdfmx
2. pdflatex footnotehyper.dtx (twice)
3. latex footnotehyper.dtx (twice), then dvips, then ps2pdf

Method (1) produces the smallest pdf files.
Options can be set in footnotehyper.tex:

- scrdoc class options (paper size, font size, ...)
- with or without source code,
- with dvipdfmx or with latex+dvips or pdflatex.

Installation:

    footnotehyper.sty    -> TDS:tex/latex/footnotehyper/footnotehyper.sty
    footnotehyper.dtx    -> TDS:source/latex/footnotehyper/footnotehyper.dtx
    footnotehyper.pdf    -> TDS:doc/latex/footnotehyper/footnotehyper.pdf
    README.md            -> TDS:doc/latex/footnotehyper/README.md

The other files may be discarded.


LICENSE
=======

This Work may be distributed and/or modified under the conditions
of the LaTeX Project Public License, version 1.3c. This version of
this license is in:

> <http://www.latex-project.org/lppl/lppl-1-3c.txt>

and the latest version of this license is in:

> <http://www.latex-project.org/lppl.txt>

Version 1.3 or later is part of all distributions of
LaTeX version 2005/12/01 or later.

The Author of this Work is:

- Jean-Francois B. `<jfbu on github>`

This Work consists of the main source file footnotehyper.dtx and the
derived files footnotehyper.sty, footnotehyper.ins, footnotehyper.tex,
footnotehyper.pdf, footnotehyper.dvi.
%</readme>
%<*tex>-------------------------------------------------------------------------
%%
%% run latex twice on this file footnotehyper.tex then dvipdfmx on
%% footnotehyper.dvi to produce the documentation footnotehyper.pdf, with
%% source code included.
%%
\chardef\Withdvipdfmx 1 % replace 1 by 0 for using latex+dvips or pdflatex
\chardef\NoSourceCode 0 % replace 0 by 1 for the doc *without* the source code
\NeedsTeXFormat{LaTeX2e}
\ProvidesFile {\pkgname.tex}[Driver for \pkgname\space documentation]%
\PassOptionsToClass   {a4paper,fontsize=11pt,oneside}{scrdoc}
\PassOptionsToPackage {english}{babel}
\input \pkgname.dtx
%%% Local Variables:
%%% mode: latex
%%% End:
%</tex>-------------------------------------------------------------------------
%<*insfile>---------------------------------------------------------------------
%%
%% Run tex on this file to extract from footnotehyper.dtx:
%%        footnotehyper.sty, footnotehyper.tex, and README.md
%%
%% Refer to README.md for installation instructions, if needed.
%%
\input docstrip.tex
\askforoverwritefalse
\def\pkgpreamble{\defaultpreamble^^J\MetaPrefix^^J%
\string\NeedsTeXFormat{LaTeX2e}^^J%
\string\ProvidesPackage{\pkgname}\perCent^^J%
\space[\pkgdate\space\pkgversion\space\pkgdescription]}%
\generate{\nopreamble\nopostamble
\file{README.md}{\from{\pkgname.dtx}{readme}}%
\usepostamble\defaultpostamble
\file{\pkgname.tex}{\from{\pkgname.dtx}{tex}}%
\usepreamble\pkgpreamble
\file{\pkgname.sty}{\from{\pkgname.dtx}{package}}}%
\catcode32=13\relax% active space
\let =\space%
\Msg{************************************************************************}
\Msg{*}
\Msg{* To finish the installation you have to move the following}
\Msg{* file into a directory searched by TeX:}
\Msg{*}
\Msg{*     \pkgname.sty} 
\Msg{*}
\Msg{* To produce the documentation run latex twice on file \pkgname.tex}
\Msg{* and then run dvipdfmx on file \pkgname.dvi.}
\Msg{*}
\Msg{* Happy TeXing!}
\Msg{*}
\Msg{************************************************************************}
\endbatchfile
%</insfile>---------------------------------------------------------------------
%<*none>------------------------------------------------------------------------
^^Afi^^Aendgroup
%
\chardef\noetex 0
\ifx\numexpr\undefined\chardef\noetex 1 \fi
\ifnum\noetex=1 \chardef\extractfiles 0 % extract files, then stop
\else
    \ifx\ProvidesFile\undefined
      \chardef\extractfiles 0 % etex etc.. on \pkgname.dtx, only file extraction.
    \else % latex/pdflatex on \pkgname.tex or on \pkgname.dtx
      \ifx\Withdvipdfmx\undefined
        % latex/pdflatex run is on \pkgname.dtx
        \chardef\extractfiles 1 % 1 = extract files and typeset manual, 2 = only typeset
        \chardef\Withdvipdfmx 0 % 0 = pdflatex or latex+dvips, 1 = dvipdfmx
        \chardef\NoSourceCode 0 % 0 =  include source code, 1 = do not
        \NeedsTeXFormat {LaTeX2e}%
        \PassOptionsToClass   {a4paper,fontsize=11pt,oneside}{scrdoc}% 
        \PassOptionsToPackage {english}{babel}%
      \else % latex run is on \pkgname.tex
        \chardef\extractfiles 2 % do not extract files, only typeset
      \fi
      \ProvidesFile{\pkgname.dtx}%
        [\pkgname\space source and documentation (\dtxtimestamp)]%
    \fi
\fi
\ifnum\extractfiles<2 % extract files
\def\MessageDeFin{\newlinechar10 \let\Msg\message
\Msg{********************************************************************^^J}%
\Msg{*^^J}%
\Msg{* To finish the installation you have to move the following^^J}%
\Msg{* file into a directory searched by TeX:^^J}%
\Msg{*^^J}%
\Msg{*\space\space\space\space \pkgname.sty^^J}%
\Msg{*^^J}%
\Msg{* To produce the documentation with source code included run latex^^J}%
\Msg{* twice on file \pkgname.tex and then dvipdfmx on \pkgname.dvi^^J}%
\Msg{*^^J}%
\Msg{* Happy TeXing!^^J}%
\Msg{*^^J}%
\Msg{********************************************************************^^J}%
}%
\begingroup
    \input docstrip.tex
    \askforoverwritefalse
    \def\pkgpreamble{\defaultpreamble^^J\MetaPrefix^^J%
    \string\NeedsTeXFormat{LaTeX2e}^^J%
    \string\ProvidesPackage{\pkgname}\perCent^^J%
    \space[\pkgdate\space\pkgversion\space\pkgdescription]}%
    \generate{\nopreamble\nopostamble
    \file{README.md}{\from{\pkgname.dtx}{readme}}%
    \usepostamble\defaultpostamble
    \file{\pkgname.ins}{\from{\pkgname.dtx}{insfile}}%
    \file{\pkgname.tex}{\from{\pkgname.dtx}{tex}}%
    \usepreamble\pkgpreamble
    \file{\pkgname.sty}{\from{\pkgname.dtx}{package}}}%
\endgroup
\fi % end of file extraction (from etex/latex/pdflatex run on \pkgname.dtx)
\ifnum\extractfiles=0 % tex/etex/xetex/etc on \pkgname.dtx, files extracted, stop
      \MessageDeFin\expandafter\end
\fi
% From this point on, run is necessarily with e-TeX.
% Check if \MessageDeFin got defined, if yes put it at end of run.
\ifdefined\MessageDeFin\AtEndDocument{\MessageDeFin}\fi
%-------------------------------------------------------------------------------
% START OF USER MANUAL TEX SOURCE
\documentclass[abstract]{scrdoc}

\ifnum\NoSourceCode=1 \OnlyDescription\fi

\usepackage{ifpdf}
\ifpdf\chardef\Withdvipdfmx 0 \fi

\makeatletter
\ifnum\Withdvipdfmx=1
   \@for\@tempa:=hyperref,bookmark,graphicx,xcolor,pict2e\do
            {\PassOptionsToPackage{dvipdfmx}\@tempa}
   %
   \PassOptionsToPackage{dvipdfm}{geometry}
   \PassOptionsToPackage{bookmarks=true}{hyperref}
   \PassOptionsToPackage{dvipdfmx-outline-open}{hyperref}
   \PassOptionsToPackage{dvipdfmx-outline-open}{bookmark}
   %
   \def\pgfsysdriver{pgfsys-dvipdfm.def}
\else
   \PassOptionsToPackage{bookmarks=true}{hyperref}
\fi
\makeatother

\usepackage[T1]{fontenc}
\usepackage[utf8]{inputenc}
\usepackage {babel}

\usepackage[hscale=0.72,vscale=0.7]{geometry}

\def\MacroFont{\ttfamily\small\hyphenchar\font45 \baselineskip11pt\relax}

\usepackage{amsmath}
\usepackage{newtxtext, newtxmath}
\usepackage[straightquotes, scaled=0.97]{newtxtt}
\usepackage{xspace}

\usepackage[dvipsnames]{xcolor}
\definecolor{joli}{RGB}{225,95,0}
\definecolor{JOLI}{RGB}{225,95,0}
\definecolor{BLUE}{RGB}{0,0,255}
\colorlet{niceone}{green!35!blue!75}

\usepackage{framed}

\usepackage[pdfencoding=pdfdoc]{hyperref}
\hypersetup{%
linktoc=all,%
breaklinks=true,%
colorlinks,%   
linkcolor=RoyalBlue,%
urlcolor=OliveGreen,%
pdfauthor={Jean-Fran\c cois B.},%
pdftitle={The \pkgname\space package},%
pdfsubject={\pkgdescription},%
pdfkeywords={LaTeX, footnotes},%
pdfstartview=FitH,%
pdfpagemode=UseNone}
% added usage of package bookmark 2013/10/10
\usepackage{bookmark} 

\usepackage{\pkgname}

\newcommand\fnh{%
        \texorpdfstring{{\color{joli}\ttfamily\bfseries \pkgname}}{\pkgname}\xspace}

\frenchspacing

\renewcommand\familydefault\sfdefault
\pagestyle{headings}

\begin{document}\thispagestyle{empty}
% the \MacroFont valid at begin document will be the one used in Implementation
% the one defined here will be used by verbatim (not by \verb)
\def\MacroFont{\ttfamily\small}
\rmfamily
\bookmark[named=FirstPage,level=1]{Title page}

\begin{center}
  {\normalfont\LARGE The \fnh  package}\\
\textsc{\small Jean-François B.}\par
  \footnotesize \ttfamily 
  jfbu on github\par
  Package version: \pkgversion\ (\pkgdate)\par
  From source file \texttt{\pkgname.dtx} of \dtxtimestamp.\par
\end{center}

\MakeShortVerb{\`}%
\begin{abstract}
The `footnote` package by Mark \textsc{Wooding} (which dates back to 1997)
allows to gather footnotes (`\savenotes`) and later insert them
(`\spewnotes`) at the bottom of the page, even if the intervening
material consists of tabulars, minipages or framed contents for
example. There is a command to patch any given environment to do
the `\savenotes/\spewnotes` commands automatically. Also there is
a `footnote` environment which allows to insert verbatim material.

The \fnh package loads `footnote.sty` and rewrites a number of its
commands to add `hyperref` awareness and to fix some issues,
particularly the one of incompatibility with `color/xcolor`. The
incompatibility with `babel-frenchb` is addressed, too. In case of
absence of `hyperref`, \fnh limits itself to partially fix some of
these issues.

This package provides no facility dedicated to footnotes from
floating environments.
\end{abstract}
\DeleteShortVerb{\`}

\section{License}

\begingroup\ttfamily\footnotesize\hyphenchar\font -1
           \parindent0pt 
           \obeyspaces\obeylines %
\pkgLicense\endgroup

\section{Usage}

DON'T load |footnote|, leave that job to \fnh. You \emph{must} load
|hyperref|.\footnote{Since |v0.9e| \fnh deactivates itself
  gracefully if |hyperref| is not loaded, or under |hyperref| option
  |hyperfootnotes=false|.}

Then you can use |\savenotes/\spewnotes| or the equivalent |savenotes|
environment; there is also a |footnote| environment, but its set-up during
\fnh loading is more delicate and it could well fail depending on how
|\@makefntext| has been customized by the class or other packages; a warning
is issued in that case. A functional |footnote| environment allows footnotes
with verbatim material.

\savenotes
\begin{framed}
Please refer to the documentation of the |footnote| package.%
\footnote{\url{http://ctan.org/pkg/footnote}}

Particularly you may check its |\makesavenoteenv| command.%
\footnote{This won't handle floating environments, though.}
\end{framed}
\spewnotes

We can try some normal footnote.\footnote{Here it is.}

\begin{savenotes}
\begin{framed}
\DeleteShortVerb{\|}\MakeShortVerb{\*}%
{\centering
  \begin{tabular}{|c|c|}
\hline
  \strut The above\footnote{Notice that if the present frame
    extended to next page,
    the end of the *savenotes* environment would then (try to) deliver its
    footnotes to that
    page.} was & coded\footnote{Alternatively a savenotes environment
    could have been used.} as:\\\hline
\end{tabular}\par}
\DeleteShortVerb{\*}\MakeShortVerb{\|}%
\begin{verbatim}
\savenotes
\begin{framed}
  Please refer to the documentation of the |footnote| package.%
  \footnote{\url{http://ctan.org/pkg/footnote}}

  Particularly you may check its |\makesavenoteenv| command.%
  \footnote{This won't handle floating environments, though.}
\end{framed}
\spewnotes
\end{verbatim}
and the present frame has \cs{footnote}'s from inside a |tabular| and is
inside a |savenotes| environment.% 
%
% plus maintenant, \small pour le \MacroFont utilisé par verbatim
%\footnote{Well, I end up on next page. Not
% much I can do about this, ask why to |framed.sty| and the output routine I
% guess.}
%
\footnote{Here is an issue which has nothing (as I finally figured out) to do
  with |footnote|, and only indirectly with \LaTeX: if you embed a
  \emph{full-width} minipage (with initial \cs{noindent}) in any environment
  not doing \cs{ignorespacesafterend}, be careful to add a \% after the
  minipage (or after the surrounding environment; or a \cs{par} immediately),
  else the output will have an extra blank line if the source has itself a
  blank line there. I hesitated adding that to
  \cs{spewnotes}/\cs{endsavenotes}. Finally I left as is.} Let's test an
|amsmath| environment with |\intertext|. As
\begin{align}
  E&=mc^2\;,
\intertext{was too easy,\footnote{There is also $E=h\nu$.}, let's
  try:}
  F&=nd^3\;.
\end{align}
\end{framed}
\end{savenotes}

Use of the \texttt{footnote} \emph{environment} with some verbatim material%
\begin{footnote}
  \verb|&$^%\[}$|
\end{footnote}
which was coded as:
\begin{verbatim}
Use of the \texttt{footnote} \emph{environment} with some verbatim material%
\begin{footnote}
  \verb|&$^%\[}$|
\end{footnote}
which was coded as:
<let's avoid infinite recursion ;-) >
\end{verbatim}
And here a use of |\footnotemark|\footnotemark\ 
followed by a |footnotetext| environment.
\begin{footnotetext}
  This one uses the normal footnote counter and the hyperlink works.
\end{footnotetext}
Now a use of |\footnotemark[N]|\footnotemark[99] followed with a
|footnotetext| environment using the same optional argument |[N]|.
No hyperfootnote link is then inserted by |hyperref|.%
\begin{footnotetext}[99]
  |hyperref| puts no hyperlink in this case, which is similar to
  the |\footnotemark[N]/\footnotetext[N]{<foo>}| context (there is
  no such problem with using the |[N]| with a |\footnote| or with
  a |footnote| environment.) 
\end{footnotetext}
And a final footnote, with a |footnote| environment using an
optional argument |[N]|.%
\begin{footnote}[100]\csname @minipagetrue\endcsname % to reduce
                                % framed inserted vertical spacing
  \begin{framed}
    \noindent\fnh deactivates itself if |hyperfootnotes=false| option to
    |hyperref| is detected, or if |hyperref| is not loaded at all.
    Essentially, it only then fixes some of the incompatibilities
    between |footnote.sty| and |color/xcolor/babel-frenchb|
    packages. But the bug making the |footnote| environment
    optional argument unusable remains unfixed by \fnh if
    |hyperref| has not been detected.
  \end{framed}
\end{footnote}
No problem with the hyperlink, then.

\section{Notes}

The |footnote| environment of |footnote.sty| had various issues
but as \fnh initially focused only on |\savenotes/\spewnotes| I
became aware only at the time of |v0.9f| of the error when this
environment is used with the optional |[N]| parameter; this is a
different bug than the one of incompatibility with |color/xcolor|.
The |v0.9f| release fixes this new problem (no extensive
tests were done.)

During this I realized (this is unrelated to |footnote| package or
\fnh) that there is a typesetting interference of |hyperref| with
|babel+frenchb| in the case of the |\footnotemark[N]| syntax. \fnh
does not fix that. I reported the problem and perhaps an upstream
fix will come at a later date (tested with |hyperref 6.83q| and
|frenchb 3.2c|).

The |hyperref| package inserts no hyperlink in case of
|\footnotemark[N]/\footnotetext[N]|. This is not modified by \fnh.

The |footnote.sty| package patches the \LaTeX\ kernel |\parbox|. I had not
realized that earlier, and starting with \fnh |v0.9f| the patch is undone.
But the re-defined macro is made available as |\fnparbox|.

\medskip

At this stage, \fnh rewrites large parts of the original
|footnote.sty| code. Perhaps a later release will not load at all
the original package (but it will keep the |\savenotes/\spewnotes|
interface.)

\StopEventually{\end{document}\endinput}
\makeatletter
    \let\check@percent\original@check@percent
\makeatother

\newgeometry{left=4cm, hscale=0.72}
\section{Implementation}

\small

\makeatletter
\noindent
\begingroup
\topsep\MacrocodeTopsep
\trivlist\parskip\z@\item[]
\macro@font
\leftskip\@totalleftmargin  \advance\leftskip\MacroIndent
\rightskip\z@  \parindent\z@  \parfillskip\@flushglue
\global\@newlistfalse \global\@minipagefalse
\ifcodeline@index
  \everypar{\global\advance\c@CodelineNo\@ne
  \llap{\theCodelineNo\ \hskip\@totalleftmargin}}%
\fi
\string\NeedsTeXFormat\string{LaTeX2e\string}\par
\string\ProvidesPackage\string{\pkgname\string}\@percentchar\par
\noindent\space [\pkgdate\space\pkgversion\space\pkgdescription]\par
\nointerlineskip
\global\@inlabelfalse
\endtrivlist
\endgroup
\makeatother

% The catcode hackery next is to avoid to have <*package> to be listed
% in the commented source code...
% (c) 2012/11/19 jf b. ;-)

\MakePercentIgnore

%
% \catcode`\<=0 \catcode`\>=11 \catcode`\*=11 \catcode`\/=11
% \let</none>\relax
% \def<*package>{\catcode`\<=12 \catcode`\>=12 \catcode`\*=12 \catcode`\/=12}
%
%</none>
%<*package>
% It is surely difficult to understand the code here as its expansion depends on
% some macros whose definitions are still in |footnote.sty| and used with no
% modification here: the most important are |\savenotes|, |\fn@startnote|, and
% |\fn@startfntext| (which is used only for the environment form.) Obviously the
% macros here depend on what these three do. 
%
% Surely, next release of \fnh will drop entirely the idea of patching
% |footnote.sty| and will contain all definitions.
% \begin{macro}{no options}
%    \begin{macrocode}
\DeclareOption*{\PackageWarning{footnotehyper}{Option `\CurrentOption' is unknown}}%
\ProcessOptions\relax
%    \end{macrocode}
% \end{macro}
% \begin{macro}{\@makefntext}
% \begin{macro}{\footnote}
% \begin{macro}{\footnotetext}
% \begin{macro}{\parbox}
% \begin{macro}{\fn@endnote}
% We load |footnote| but leave to the user to take care of |hyperref|.
%
% As we need to intercept some re-definitions done by
% |footnote.sty| we will first check if it is already loaded.
%
% If \cs{@makefntext} has been customized and its argument is not visible at
% top level in its meaning, then loading |footnote.sty| will fail with a low
% level \TeX\ error. We save its meaning and replace it by an innocuous one
% for the time being. We will come back to this at begin document.
%
% |footnote.sty| modifies |\parbox|. We (|v0.9f|) undo the patch, because
% modifying |\parbox| in any way for no good reason is BAD.
%
% The |\fn@endnote| was |\let| to |\color@endgroup| which is bad as it fails if
% |color/xcolor| is not yet loaded. Besides its presence was missing from some
% needed place (see later).
% 
% We postpone to at begin document the redefinitions of \cs{footnote} and
% \cs{footnotetext}.
%    \begin{macrocode}
\@ifpackageloaded{footnote}
 {\PackageWarningNoLine{footnotehyper}%
  {Please next time do not load package footnote,^^J
 but leave that to me. For the time being, I will start by undoing the^^J
 modifications it did to \string\parbox. Use \string\fnparbox\space if needed.}%
 }
 {\let\FNH@@makefntext\@makefntext\let\@makefntext\@firstofone
  \RequirePackage{footnote}
  \let\@makefntext\FNH@@makefntext
 }%
\let\fnparbox\parbox\let\parbox\fn@parbox
\let\FNH@fn@footnote    \footnote     % footnote.sty's \footnote will get
\let\FNH@fn@footnotetext\footnotetext % redefined later. We preserve also
\let\footnote    \fn@latex@@footnote  % the meaning of \footnote from prior
\let\footnotetext\fn@latex@@footnotetext % to the loading of footnote.sty
\def\fn@endnote  {\color@endgroup}%
%    \end{macrocode}
% \end{macro}\end{macro}\end{macro}\end{macro}\end{macro}
% There are some |\let|'s done by |footnote.sty| in what appears to be
% premature ways.
%    \begin{macrocode}
\AtBeginDocument {%
    \let\fn@latex@@footnote    \footnote % meaning of \footnote at end of preamble
    \let\fn@latex@@footnotetext\footnotetext
}%
%    \end{macrocode}
% \begin{macro}{\FNH@inactive}
% Apart from fixing some of the bugs with |footnote.sty|'s footnote
% environment, we intervene only if |hyperfootnotes| option of |hyperref|
% applies.
%    \begin{macrocode}
\AtBeginDocument{\@ifpackageloaded{hyperref}
  {\ifHy@hyperfootnotes
    \let\fn@fntext \FNH@hyper@fntext
    \let\spewnotes \FNH@hyper@spewnotes
    \let\endsavenotes\spewnotes
    \let\fn@endfntext\FNH@fixed@endfntext
    \let\footnote    \FNH@fixed@footnote
    \let\footnotetext\FNH@fixed@footnotetext
   \else
     \FNH@inactive
   \fi }{\FNH@inactive}% 
    \let\endfootnote\fn@endfntext
    \let\endfootnotetext\endfootnote
}%
%    \end{macrocode}
% |v0.9e| did a partial fix of the |footnote| environment from |footnote.sty|,
% |v0.9f| fix is more complete but is applied only if |hyperref| is present with
% option |hyperfootnotes=true|.
%    \begin{macrocode}
\def\FNH@fix@endfntext\@finalstrut\strutbox\fn@postfntext
    {\@finalstrut\strutbox\fn@postfntext\fn@endnote}%
\def\FNH@inactive {%
    \let\footnote    \FNH@fn@footnote
    \let\footnotetext\FNH@fn@footnotetext
    \expandafter\expandafter\expandafter\def
    \expandafter\expandafter\expandafter\fn@endfntext
    \expandafter\expandafter\expandafter
       {\expandafter\FNH@fix@endfntext\fn@endfntext}%
    \PackageInfo{footnotehyper}%
     {hyperref package not loaded or hyperfootnotes^^J
  =false option; I did not activate myself and only patched footnote.sty^^J
  for color/xcolor compatibility, leaving some other issues unfixed.}}%
%    \end{macrocode}
% \end{macro}
% \begin{macro}{\FNH@hyper@fntext}
% Some |amsmath| complications. I change the coding, but same effect as
% in original |\fn@fntext|.
%    \begin{macrocode}
\def\FNH@hyper@fntext {\ifx\ifmeasuring@\undefined\expandafter\@secondoftwo
                                      \else\expandafter\@firstofone\fi
    {\ifmeasuring@\expandafter\@gobbletwo\fi}%
    \FNH@hyper@fntext@i }%
%    \end{macrocode}
% \end{macro}
% \begin{macro}{\FNH@hyper@fntext@i}
% Note: original \cs{fn@fntext} had no \cs{long}, which looks wrong.
%
% We do the \cs{ifHy@nesting} test although hyperref's manual
% says ``Allows links to be nested; no drivers currently support this.''
%    \begin{macrocode}
\long\def\FNH@hyper@fntext@i #1{\global\setbox\fn@notes\vbox
     {\unvbox\fn@notes
      \fn@startnote
      \@makefntext
       {\rule\z@\footnotesep\ignorespaces
        \ifHy@nesting\expandafter\ltx@firstoftwo
                \else\expandafter\ltx@secondoftwo
        \fi
        {\expandafter\hyper@@anchor\expandafter{\Hy@footnote@currentHref}{#1}}%
        {\Hy@raisedlink
          {\expandafter\hyper@@anchor\expandafter{\Hy@footnote@currentHref}%
          {\relax}}%
         \let\@currentHref\Hy@footnote@currentHref
         \let\@currentlabelname\@empty
         #1}%
      \@finalstrut\strutbox }%
      \fn@endnote }%
}%
%    \end{macrocode}
% \end{macro}
% \begin{macro}{\FNH@hyper@spewnotes}
%   The final touch in our hack is to patch the original \cs{spewnotes} for it
%   to use the original, non-hyperref modified, version of \cs{@footnotetext}.
%   And let's not forget \cs{endsavenotes} (done at begin document).
%    \begin{macrocode}
\def\FNH@hyper@spewnotes {\endgroup
    \if@savingnotes\else\ifvoid\fn@notes\else
    \begingroup\let\@makefntext\@empty
               \let\@finalstrut\@gobble
               \let\rule\@gobbletwo
               \H@@footnotetext{\unvbox\fn@notes}%
    \endgroup\fi\fi
}%
%    \end{macrocode}
% \end{macro}
% \begin{macro}{\FNH@fixed@endfntext}
%   We now take care of |footnote.sty|'s |footnote| environment. The original
%   \cs{fn@endfntext} is lacking a \cs{fn@endnote}, and this meant that
%   |footnote.sty| was incompatible with |color/xcolor| packages. Also this
%   \cs{fn@endnote} was |\let| to |\color@endgroup| which is wrong.
%
%   Furthermore, independently of presence of the |\color/xcolor| issue, the
%   |footnote.sty|'s |footnote| environment raised an error if used with an
%   optional argument. |v0.9f| addresses this issue.
%
%   The |footnotetext| environment adds a complication, in case of optional
%   argument we should not try to set up a link due to the fact that |hyperref|
%   does not support them for |\footnotemark[N]/\footnotetext[N]| syntax.
%   And we need to make sure that the |footnote| and |footnotetext| environments
%   obey the |\savenotes/\spewnotes| mechanism.
%
%   To handle all of this we code things completely differently from
%   |footnote.sty|, its |\fn@getmark| isn't used at all.
%    \begin{macrocode}
\def\FNH@fixed@endfntext {%
    \@finalstrut\strutbox
    \fn@postfntext
    \fn@endnote
    \egroup\FNH@endfntext@next % will decide if link or no link
}%
\def\FNH@endfntext@link {\begingroup
    \let\@makefntext\@empty\let\@finalstrut\@gobble\let\rule\@gobbletwo
    \@footnotetext {\unvbox\z@}%
    \endgroup
}%
\def\FNH@endfntext@nolink {\begingroup
    \let\@makefntext\@empty\let\@finalstrut\@gobble
    \let\rule\@gobbletwo
    \if@savingnotes\expandafter\fn@fntext\else\expandafter\H@@footnotetext\fi
    {\unvbox\z@}\endgroup
}%
%    \end{macrocode}
% \end{macro}
% \begin{macro}{\FNH@fixed@footnote}
% \begin{macro}{\FNH@fixed@footnotetext}
%    \begin{macrocode}
\def\FNH@fixed@footnote {\ifx\@currenvir\fn@footnote
    \expandafter\FNH@footnoteenv\else\expandafter\fn@latex@@footnote\fi }%
\def\FNH@footnoteenv {\@ifnextchar[\FNH@xfootnoteenv%]
    {\stepcounter\@mpfn
     \protected@xdef\@thefnmark{\thempfn}\@footnotemark
     \def\FNH@endfntext@next{\FNH@endfntext@link}\fn@startfntext}}%
\def\FNH@xfootnoteenv [#1]{%
    \begingroup
     \csname c@\@mpfn\endcsname #1\relax
     \unrestored@protected@xdef\@thefnmark{\thempfn}%
    \endgroup\@footnotemark\def\FNH@endfntext@next{\FNH@endfntext@link}%
    \fn@startfntext}%
\def\FNH@fixed@footnotetext {\ifx\@currenvir\fn@footnotetext
  \expandafter\FNH@footnotetextenv\else\expandafter\fn@latex@@footnotetext\fi}%
\def\FNH@footnotetextenv {\@ifnextchar[\FNH@xfootnotetextenv%]
    {\protected@xdef\@thefnmark{\thempfn}%
     \def\FNH@endfntext@next{\FNH@endfntext@link}\fn@startfntext}}%
\def\FNH@xfootnotetextenv [#1]{%
    \begingroup
     \csname c@\@mpfn\endcsname #1\relax
     \unrestored@protected@xdef\@thefnmark{\thempfn}%
    \endgroup\def\FNH@endfntext@next{\FNH@endfntext@nolink}%
    \fn@startfntext }%
%    \end{macrocode}
% \end{macro}\end{macro}
% \begin{macro}{\@makefntext}
% \begin{macro}{\fn@prefntext}
% \begin{macro}{\fn@postfntext}
%   The definitions of |\fn@prefntext| and |\fn@prefntext| (which are needed for
%   the |footnote| environment, within |\fn@startfntext| and |\fn@endfntext|)
%   are from a daring analysis done inside |footnote.sty| of |\@makefntext|. It
%   may very easily fail to give a sensible result. \fnh does a check of its own
%   to determine if that was the case but can't salvage the situation, except
%   for hard-coded rescues like done here with |frenchb|. In particular the
%   mechanism is doomed to failure if the parameter in |\@makefntext| is not
%   visible at top level in its meaning. It is also wrong if the parameter is
%   used multiple times. These checks will fail (rather, they are skipped) if
%   |footnote|'s loading was done by the user and not by |footnotehyper|.
%
%   Note: there is an interference which I could fix here between |hyperref|
%   and |frenchb| regarding |\footnotemark[NUM]|, but it is better to wait it
%   is fixed on |frenchb| side (or |hyperref|'s side !), else my own fix might
%   badly interfere with the ``official'' one |;-)|.
%    \begin{macrocode}
\ifx\FNH@@makefntext\undefined\expandafter\@gobble
   \else\expandafter\AtBeginDocument\fi
{%
 \ifx\@makefntextFB\undefined
                   \expandafter\@gobble\else\expandafter\@firstofone\fi
 {\ifFBFrenchFootnotes \let\FNH@@makefntext\@makefntextFB \else 
                       \let\FNH@@makefntext\@makefntextORI\fi}% 
 \expandafter\FNH@check@a\FNH@@makefntext{1.2!3?4,}\FNH@@@1.2!3?4,\FNH@@@\relax
}%
\long\def\FNH@check@a #11.2!3?4,#2\FNH@@@#3%
{%
     \ifx\relax#3\expandafter\@firstoftwo\else\expandafter\@secondoftwo\fi
     \FNH@bad@footnote@env
     {\def\fn@prefntext{#1}\def\fn@postfntext{#2}\FNH@check@b}%
}%
\def\FNH@check@b #1\relax
{%
    \expandafter\expandafter\expandafter\FNH@check@c
    \expandafter\meaning\expandafter\fn@prefntext
                            \meaning\fn@postfntext1.2!3?4,\FNH@check@c\relax
}%
\def\FNH@check@c #11.2!3?4,#2#3\relax 
   {\ifx\FNH@check@c#2\expandafter\@gobble\fi\FNH@bad@footnote@env}%
\def\FNH@bad@footnote@env 
{\PackageWarningNoLine{footnotehyper}%
 {The footnote environment from package footnote^^J
  will be dysfunctional, sorry (not my fault...). You may try to mail me the^^J
  preamble and/or only the next lines:}%
    \typeout{\meaning\@makefntext}%
    \let\fn@prefntext\@empty\let\fn@postfntext\@empty  
}%
\endinput
%    \end{macrocode}
% \end{macro}\end{macro}\end{macro}
% \MakePercentComment
\Finale
%%
%% End of file `footnotehyper.dtx'.
